\begin{frame}[fragile]{실습: 4$\times$4 특징지도를 최대풀링으로}
  10.1절의 \texttt{conv2d} 코드와 같은 스타일로 최대풀링 직접 구현:
  \begin{columns}[T]
    \begin{column}{0.60\textwidth}
\begin{lstlisting}
def max_pool(feature_map, pool_size=2):
    h, w = len(feature_map), \
           len(feature_map[0])
    out_h, out_w = h // pool_size, \
                   w // pool_size
    output = [[0.0] * out_w for _
              in range(out_h)]
    for i in range(out_h):
        for j in range(out_w):
            block = [
                feature_map[
                    i*pool_size+di]
                    [j*pool_size+dj]
                for di in range(pool_size)
                for dj in range(pool_size)]
            output[i][j] = max(block)
    return output

fm = [[1, 0, 2, 1],
      [1, 0, 3, 1],
      [1, 0, 4, 1],
      [1, 0, 2, 1]]
print(max_pool(fm, 2))  # [[1, 3], [1, 4]]
\end{lstlisting}
    \end{column}
    \begin{column}{0.40\textwidth}
      \begin{itemize}
        \item ``원본'' 특징지도를 그대로 넣으면
          $\begin{bmatrix}1&3\\1&4\end{bmatrix}$ ---
          앞 계산과 \textbf{정확히 일치}
        \item 1픽셀 이동 버전도 \textbf{같은 결과}
          --- 둔감함이 코드로
      \end{itemize}
    \end{column}
  \end{columns}
\end{frame}
